\begin{tikzpicture}[
    >=Latex,
    font=\small,
    nota/.style={font=\footnotesize, align=center},
    flujo/.style={-{Latex[width=5pt,length=4pt]}, thick, gray}
]

% Tolva genérica de salida inferior: centro de la salida en (#1,#2)
\newcommand{\tolvasup}[2]{%
    \draw[thick] (#1-1.0,#2+1.6) -- (#1+1.0,#2+1.6) -- (#1+0.3,#2+0.4) --
        (#1+0.3,#2) (#1-0.3,#2) -- (#1-0.3,#2+0.4) -- (#1-1.0,#2+1.6);
}

% (a) Tornillo con agitador
\tolvasup{0}{2.0}
\draw[thick] (0,3.5) -- (0,2.4);
\foreach \y in {2.7,3.0,3.3}{\draw (-0.25,\y) -- (0.25,\y);}
\draw[gray] (0.25,3.1) -- (1.2,3.45);
\node[nota, anchor=west] at (1.2,3.45) {Agitador};
\draw[thick] (-0.3,1.6) rectangle (1.9,2.0);
\foreach \x in {-0.1,0.3,0.7,1.1,1.5}{\draw (\x,1.6) -- (\x+0.25,2.0);}
\draw[thick, fill=gray!25] (-0.85,1.6) rectangle (-0.3,2.0);
\node[nota] at (-0.6,1.25) {Motor};
\draw[flujo] (1.75,1.55) -- (1.75,0.9);
\node[nota] at (0.8,0.45) {(a) Tornillo con agitador};

% (b) Canal vibratorio
\tolvasup{4.6}{2.0}
\draw[thick] (4.1,1.95) -- (6.3,1.65);
\draw[thick] (4.1,1.95) -- (4.1,2.05);
\draw[thick, fill=gray!25] (4.6,1.25) rectangle (5.4,1.6);
\node[nota] at (5.0,1.05) {Excitador};
\draw (5.0,1.6) -- (4.94,1.64) -- (5.06,1.69) -- (4.94,1.74) -- (5.0,1.78);
\draw[flujo] (6.2,1.6) -- (6.2,0.9);
\node[nota] at (5.2,0.45) {(b) Canal vibratorio};

% (c) Compuerta
\tolvasup{9.0}{1.8}
\draw[very thick] (8.7,1.75) -- (9.15,1.55);
\fill (8.7,1.75) circle (1.5pt);
\node[nota, anchor=west] at (9.4,1.55) {Compuerta con\\servomotor};
\draw[flujo] (9.0,1.45) -- (9.0,0.9);
\node[nota] at (9.0,0.45) {(c) Compuerta};

\end{tikzpicture}
